% Repertorio bibliografico unico del volume I, aggiornato per il capitolo su Talete.
% Per le fonti antiche indichiamo i passi stabili; le edizioni e traduzioni non
% ancora selezionate saranno registrate quando effettivamente utilizzate.
\chapter{Fonti e bibliografia}

\section*{Fonti antiche}

\begin{description}
  \item[Erodoto] \emph{Storie}, I, 74--75 e 170.
  \item[Platone] \emph{Protagora}, 343a; \emph{Teeteto}, 174a.
  \item[Aristotele] \emph{Metafisica}, I, 3, 983b6--27 (DK 11 A12);
    \emph{Politica}, I, 11, 1259a10--16 (DK 11 A10);
    \emph{De caelo}, II, 13, 294a28--34;
    \emph{De anima}, I, 2, 405a19--21 e I, 5, 411a7--8.
  \item[Seneca] \emph{Questioni naturali}, III, 14 (DK 11 A15).
  \item[Diogene Laerzio] \emph{Vite dei filosofi}, I, 22--44, in particolare
    I, 23, 27 e 37--38 (DK 11 A1).
  \item[Proclo] \emph{Commento al primo libro degli Elementi di Euclide},
    p.~65, ed.\ G.~Friedlein, Leipzig, Teubner, 1873
    (testimonianza sulla tradizione geometrica risalente a Eudemo di Rodi).
\end{description}

\section*{Raccolte di testimonianze}

Hermann Diels e Walter Kranz (a cura di), \emph{Die Fragmente der Vorsokratiker},
6\textsuperscript{a} ed., Berlin, Weidmann, 1951--1952, cap.~11 (Talete).
Le sigle DK utilizzate nel testo rinviano alla numerazione di questa raccolta.

\section*{Studi moderni}

Lidia Palumbo, \emph{Trentadue ore di filosofia antica},
2\textsuperscript{a} ed., Napoli, Loffredo, 2015,
«Prima ora», in particolare pp.~17--19.

Emanuele Severino, \emph{La filosofia dai Greci al nostro tempo},
vol.~I, \emph{La filosofia antica e medioevale},
Milano, BUR, 2010, «Talete», pp.~48--50.

Friedrich Nietzsche, \emph{La filosofia nell'epoca tragica dei Greci}.
L'interpretazione discussa nel capitolo è ripresa attraverso
Palumbo, \emph{Trentadue ore}, p.~19; non è stata consultata
direttamente una specifica edizione di quest'opera.
